% Dades pròpies del llibre. Es carrega abans del preàmbul compartit.
% Una sola font de veritat per a la portada i les metadades del PDF.
% Actualitzeu versió i data a cada versió publicada.
\newcommand{\llibretitol}{Teoria de conjunts}
\newcommand{\llibresubtitol}{Una visió categòrica sense categories}
\newcommand{\llibreautor}{Miquel Àngel Perelló}
\newcommand{\llibreversio}{1.7.1}
% Data editorial fixa (reproduïble), en lloc de \today.
\newcommand{\llibredata}{Octubre de 2026}
\newcommand{\llibreany}{2026}
% Els enllaços ja inclouen \url (no s'han d'embolcallar de nou).
\newcommand{\llibreweb}{\url{https://aprendes.com/logcat/}}
\newcommand{\llibrerepositori}{\url{https://github.com/maperello/aprendes}}
\newcommand{\llibreerrates}{\url{https://github.com/maperello/aprendes/issues}}

\title{\llibretitol\\[0.8em]\large\llibresubtitol}
\author{\llibreautor}
\date{\llibredata}

% S'ha de cridar després de carregar hyperref (és a preamble.tex).
\newcommand{\hypermetadata}{%
  \hypersetup{%
    pdftitle={\llibretitol: \llibresubtitol},
    pdfauthor={\llibreautor},
    pdfsubject={Teoria estructural de conjunts (ETCS) exposada en un llenguatge de primer ordre sobre aplicacions, sense pressuposar la teoria de categories: parts, construccions universals, exponencials, nombres naturals, cardinalitat, bons ordres, ordinals i cardinals. Versió \llibreversio\ (\llibredata). Llicència CC BY-SA 4.0},
    pdfkeywords={teoria de conjunts, ETCS, conjunts estructurals, Lawvere, topos ben apuntat, axioma d'elecció, ordinals, cardinals, hipòtesi del continu},
    bookmarksnumbered=true,
    bookmarksopen=true%
  }%
}
